\section{From the height profile to the spin mass}\label{sec:mass}

The height profile identifies a transition at the logarithmic scale
\(T\).  We now show that the correlation length has this same scale.
Positive pinned-height covariance below \(T\) prevents a substantially
shorter correlation length.  Above \(T\), small free-height fluctuations
allow a circulation estimate with an arbitrarily large fixed tilt;
a finite-size correlation inequality then gives the opposite bound.
Both arguments concern the free-box thermodynamic correlation used in
Theorem~\ref{thm:main}.

Throughout this section, fix a sequence \(b_k< b_c\) tending to \(b_c\),
put \(\delta_k=\sqrt{b_c-b_k}\), and pass to a subsequence furnished by
Proposition~\ref{prop:profile-step}.  On this subsequence the height
profile satisfies
\begin{equation}\label{eq:mass-profile}
 \mathcal A(w)=a\quad(0<w<T),\qquad
 \mathcal A(w)=0\quad(w>T),\qquad a=8\pi.
\end{equation}
All finite-graph estimates below are uniform when \(b\) belongs to a
fixed compact interval \(I\subset(0,\infty)\) containing the tail of
\((b_k)\).  We will prove
\begin{equation}\label{eq:mass-scale-goal}
             \delta_k\log_L\frac1{m(b_k)}\longrightarrow T.
\end{equation}

\subsection{Correlation inequalities with the prescribed mass}

For a finite graph \(G\), write
\[
 C_{G,b}(u,v)=\left\langle e^{i(\theta_u-\theta_v)}\right\rangle_{G,b},
 \qquad
 C^{(2)}_{G,b}(u,v)=
 \left\langle e^{2i(\theta_u-\theta_v)}\right\rangle_{G,b}.
\]
Here every edge has the cosine interaction with parameter \(b\), and
the boundary is free.  Reflection of all angles makes these quantities
real, and their Fourier expansions make them nonnegative.  We use
\(C_b(x)\) for the infinite-volume free correlation between \(0\) and
\(x\in\Z^2\), so \(C_b(r)=C_b(re_1)\).

\begin{lemma}\label{lem:mass-correlation-bounds}
For every \(b>0\) and \(x\in\Z^2\),
\begin{equation}\label{eq:mass-pointwise}
                       C_b(x)\le e^{-m(b)\|x\|_\infty}.
\end{equation}
If \(G\) is a finite subgraph of \(\Z^2\), then
\begin{equation}\label{eq:mass-charge-two}
 0\le C^{(2)}_{G,b}(u,v)\le4C_{G,b}(u,v)
                         \le4C_b(v-u).
\end{equation}
\end{lemma}

\begin{proof}
The Ginibre inequalities \cite{Ginibre} imply positive correlation
for cosine characters and monotonicity of two-point correlations
under increasing ferromagnetic edge couplings.  In particular,
finite free correlations are bounded above by their free-exhaustion
limits; see also \cite[Lemma~16]{vEL}.  These limits are translation
invariant: translating any finite graph gives another finite graph
in the same monotone exhaustion.

We first record supermultiplicativity.  For three vertices \(u,v,z\),
Ginibre's inequality gives
\[
 \left\langle
 \cos(\theta_u-\theta_v)\cos(\theta_z-\theta_v)
 \right\rangle
 \ge C_{G,b}(u,v)C_{G,b}(z,v).
\]
The expectation of
\(\sin(\theta_u-\theta_v)\sin(\theta_z-\theta_v)\) is nonnegative.
To see this, condition on \(\theta_v\), rotate it to zero, and then
condition on all cosine coordinates.  Writing a sine coordinate
as its absolute value times its sign, the remaining signs have a
ferromagnetic zero-field Ising distribution: the sign coupling on
an edge \(\{x,y\}\) is
\(b|\sin\theta_x\sin\theta_y|\ge0\).
The sign at \(v\), whose sine is zero, is irrelevant.
Their two-point sign expectations are therefore nonnegative.
The cosine addition formula now proves
\[
               C_{G,b}(u,z)\ge C_{G,b}(u,v)C_{G,b}(v,z).
\]
After free exhaustion this becomes
\(C_b(x+y)\ge C_b(x)C_b(y)\).
For every positive integer \(n\),
\(C_b(jne_1)\ge C_b(ne_1)^j\).  Taking \(j\to\infty\) in the
definition of the axis mass yields
\(C_b(ne_1)\le e^{-m(b)n}\).
If \(x=(p,q)\) and \(|p|\ge|q|\), reflection symmetry and
supermultiplicativity give
\[
 C_b(x)^2=C_b(p,q)C_b(p,-q)\le C_b(2p,0)
                                      \le e^{-2m(b)|p|}.
\]
Interchanging coordinates treats the other case and proves
\eqref{eq:mass-pointwise} with the mass specified in the statement
of the main theorem.

For the charge-two estimate, fix an axis and abbreviate
\(s_x=\sin\theta_x\), \(c_x=\cos\theta_x\).
Rotation invariance gives
\[
 C_{G,b}(u,v)=2\langle s_us_v\rangle_{G,b},\qquad
 C^{(2)}_{G,b}(u,v)=8\langle c_uc_vs_us_v\rangle_{G,b}.
\]
The same conditional sign argument gives
\(\langle s_us_v\mid(c_x)_{x\in G}\rangle\ge0\).
Since \(c_uc_v\le1\), the second identity is at most
\(8\langle s_us_v\rangle_{G,b}=4C_{G,b}(u,v)\).
Edge monotonicity finishes \eqref{eq:mass-charge-two}.
\end{proof}

\subsection{Pinned covariance gives a lower bound on length}

We use the random-loop representation of the XY current due to
van Engelenburg and Lis \cite[Section~3]{vEL}.  We describe its
normalization because the height boundary condition and the loop
orientation are important here.
For each oriented edge of a finite graph \(G\), let
\(n_{uv}\in\Z_{\ge0}\) be a current.  A current with zero divergence at
every vertex has weight
\begin{equation}\label{eq:mass-current-weight}
            \prod_{\{u,v\}\in E(G)}
       \frac{(b/2)^{n_{uv}+n_{vu}}}{n_{uv}!\,n_{vu}!}.
\end{equation}
Equivalently, take independent Poisson variables of mean \(b/2\)
on the oriented edges and condition on zero divergence.
For an unoriented edge put \(N_e=n_{uv}+n_{vu}\), label its
\(N_e\) copies, and assign an orientation to each copy.
Each labeled assignment receives edge weight
\((b/2)^{N_e}/N_e!\).
At a vertex with \(\ell_v\) incoming and \(\ell_v\) outgoing
half-edges, pair them by a uniformly chosen bijection.
The resulting directed loops have weight
\begin{equation}\label{eq:mass-loop-weight}
 \prod_{e\in E(G)}\frac{(b/2)^{N_e}}{N_e!}
                     \prod_{v\in V(G)}\frac1{\ell_v!}.
\end{equation}
Indeed, summing orientations introduces the binomial factor
\(N_e!/(n_{uv}!n_{vu}!)\), and summing pairings cancels the vertex
factor.  Thus this is a probability extension of the normalized
current law \eqref{eq:mass-current-weight}.

Let \(\mathbf P_{G,b}\) denote this loop law.  Given the undirected
paired loops, their orientations are independent and fair.
The reason is that \(\ell_v=\deg(v)/2\) is fixed by that paired
geometry, so every choice of loop orientations has the same weight
in \eqref{eq:mass-loop-weight}.  For a planar rectangle, current
integration with exterior face height zero gives
\begin{equation}\label{eq:mass-winding-height}
                  h(z)=2\pi\sum_\gamma W_\gamma(z),
\end{equation}
where \(W_\gamma(z)\) is the oriented winding number of the loop
\(\gamma\) around the bounded face \(z\)
\cite[Proposition~7]{vEL}.  The factor \(2\pi\) accounts for our
height spacing.

We will also use the following consequence of the loop switching
identity \cite[Lemma~8]{vEL}:
\begin{equation}\label{eq:mass-loop-connection}
 \mathbf P_{G,b}\{\text{a loop visits both }u\text{ and }v\}
                         \le2C^{(2)}_{G,b}(u,v).
\end{equation}
In the notation of that lemma, let \(m_{uv}\) be the number of
directed pieces of loops from \(u\) to \(v\) with no internal visit
to either endpoint.  A loop visiting both endpoints implies
\(m_{uv}\ge1\), and the identity is
\(C^{(2)}_{G,b}(u,v)=\mathbf E_{G,b}[m_{uv}/(m_{uv}+1)]\).
Thus \eqref{eq:mass-loop-connection} applies to precisely the
normalized loop law above.

\begin{lemma}[Separated height covariance]\label{lem:mass-loop-covariance}
Fix \(d,M,c_1>0\) and a compact \(I\subset(0,\infty)\).
Let \(G_n\) be a free spin rectangle contained in a square of side
\(c_1n\).  On its bounded faces take the dual height law with
exterior face height zero.  Let
\[
               F_n=\sum_z f_{n,z}h(z),\qquad
               H_n=\sum_z g_{n,z}h(z),
\]
where \(\sum_z|f_{n,z}|,\sum_z|g_{n,z}|\le M\), and the two
supports are separated horizontally by at least \(dn\).
There are constants \(c,C>0\), depending only on the displayed
fixed parameters, such that, for \(b\in I\) and sufficiently
large \(n\),
\begin{equation}\label{eq:mass-covariance-decay}
 |\Cov(F_n,H_n)|\le
       Cn^{12}e^{-c n m(b)}+Ce^{-c n^4\log n}.
\end{equation}
\end{lemma}

\begin{proof}
Condition on the undirected paired geometry.  Choose one reference
orientation for each loop and write its random orientation as
\(\sigma_\gamma\in\{-1,1\}\).  Equation~\eqref{eq:mass-winding-height}
gives
\[
 F_n=\sum_\gamma\sigma_\gamma F_n(\gamma),\qquad
 H_n=\sum_\gamma\sigma_\gamma H_n(\gamma),
\]
where
\(F_n(\gamma)=2\pi\sum_zf_{n,z}W_\gamma^{\rm ref}(z)\),
and similarly for \(H_n(\gamma)\).
The conditional means vanish, and the conditional covariance is
\(\sum_\gamma F_n(\gamma)H_n(\gamma)\).
If one summand is nonzero, the loop has nonzero winding around
at least one face in each support.  Each such face lies in the
horizontal range of the loop.  Consequently this loop has
horizontal diameter at least \(dn-O(1)\).

Let \(N=\sum_eN_e\) be the total number of traversals, and let
\(\mathcal B_n\) be the event that some loop has horizontal diameter
at least \(dn/2\).  A loop of length \(|\gamma|\) has winding
number at most \(|\gamma|\) in absolute value.  Hence
\[
 \left|\sum_\gamma F_n(\gamma)H_n(\gamma)\right|
 \le (2\pi M)^2\sum_\gamma|\gamma|^2\,\mathbf1_{\mathcal B_n}
 \le (2\pi M)^2N^2\mathbf1_{\mathcal B_n}.
\]
There are \(O(n^4)\) ordered pairs of vertices in \(G_n\).
A loop counted by \(\mathcal B_n\) visits a pair separated by at
least \(dn/2\).  Using \eqref{eq:mass-loop-connection} and
Lemma~\ref{lem:mass-correlation-bounds}, we obtain
\[
 \mathbf P_{G_n,b}(\mathcal B_n)
 \le C n^4 e^{-c n m(b)}.
\]
On \(N\le n^4\), the preceding conditional-covariance bound is
therefore at most \(C n^{12}e^{-c n m(b)}\) after expectation.

It remains to bound the weighted tail of \(N\).  Before the
zero-divergence conditioning, \(N\) is Poisson with mean
\(\mu=b|E(G_n)|\le C_I n^2\).  The probability of that conditioning
event is at least the probability of the zero configuration,
namely \(e^{-\mu}\).  Thus
\[
 \mathbf E_{G_n,b}[N^2\mathbf1_{\{N>n^4\}}]
 \le e^{\mu}\,
     \E_{\operatorname{Pois}(\mu)}[N^2\mathbf1_{\{N>n^4\}}]
 \le Ce^{-c n^4\log n}.
\]
For the last estimate one may use the Poisson exponential moment
with parameter \(\log n\): its logarithm is
\(\mu(n-1)=O(n^3)\), whereas the tail factor is
\(e^{-n^4\log n}\); the polynomial insertion changes only a
polynomial prefactor.  Pairings do not change the marginal law of
\(N\).  This proves \eqref{eq:mass-covariance-decay}.
\end{proof}

We apply the lemma to the nondegenerate pinned limits supplied by
Proposition~\ref{prop:gaussian-profile}.  Let \(D=(0,1)^2\), take a
larger fixed rectangle \(U\), and pin every primary height in a
positive-width rim \(p=U\setminus D\), with the usual admissible
rounding at mesh \(n^{-1}\).  Choose the unpinned vertices to be
exactly the \(n\)-by-\(n\) grid of bounded faces of the spin rectangle
\(\{0,\ldots,n\}^2\).  This is only a translation and an
\(O(n^{-1})\) rounding of the primary mesh.

The resulting interior height law is exactly the dual law of that
\emph{free} spin rectangle.  All pin--pin edge factors cancel from
the normalized law.  Every edge from an unpinned face to the pinned
rim becomes an edge to the exterior face; their multiplicities are
retained.  In particular a corner face has two such exterior edges.
There are no fixed boundary spins or additional period constraints.
This agrees with free-boundary duality
\cite[Lemma~2.1]{X}.

Fix smooth nonzero functions \(f,g\ge0\), compactly supported in two
interior boxes in \(D\) with positive horizontal separation, and
put
\[
 F_n=n^{-2}\sum_z f(z/n)h(z),\qquad
 H_n=n^{-2}\sum_z g(z/n)h(z).
\]
The harmless translations of face centers are included in the
argument of the sampled functions.  If
\(\delta_k\log_L n_k\to w\in(0,T)\),
Proposition~\ref{prop:gaussian-profile} gives
\begin{equation}\label{eq:mass-positive-pinned}
 \Cov_{b_k,p}(F_{n_k},H_{n_k})
 \longrightarrow a\int_D\int_Df(z)G_D(z,z')g(z')\,dz\,dz'>0,
\end{equation}
where \(G_D\) is the Dirichlet Green kernel on \(D\).
The strict inequality follows from its positivity in the connected
square.  These smears have nonzero total mass.  Their convergence
uses the grounded inverse-form part of
Proposition~\ref{prop:gaussian-profile}: the primary zero rim fixes
the constant, and the approximation cost is the grounded cost on
\(D\).  No mean-zero condition is imposed on \(f\) or \(g\).

\begin{proposition}\label{prop:mass-lower-length}
Along the extracted sequence,
\begin{equation}\label{eq:mass-liminf}
                 \liminf_k\delta_k\log_L\frac1{m(b_k)}\ge T.
\end{equation}
\end{proposition}

\begin{proof}
If the assertion failed, a further subsequence and numbers
\(0<w<T\), \(\eta>0\) would satisfy
\[
                  \delta_k\log_L\frac1{m(b_k)}\le w-2\eta.
\]
Take \(n_k=\lfloor L^{w/\delta_k}\rfloor\).  Then, for all large
\(k\),
\[
 n_km(b_k)\ge L^{\eta/\delta_k}
                    \ge n_k^{\eta/(2w)}.
\]
The coefficients of \(F_{n_k}\) and \(H_{n_k}\) have uniformly
bounded absolute sums.  Lemma~\ref{lem:mass-loop-covariance} makes
their covariance tend to zero, because an exponential in a positive
power of \(n_k\) dominates the polynomial in
\eqref{eq:mass-covariance-decay}.  This contradicts
\eqref{eq:mass-positive-pinned}.
\end{proof}

\subsection{Zero coefficient on an interval gives strong decay}

We next use the part of the profile above \(T\).  The argument is a
finite-scale version of the circulation proof of
\cite[Proposition~2.6]{X}.  Its relevant input is especially simple:
on a compact interval of logarithmic scales with zero coefficient,
the transforms of bounded affine tests on free squares tend
uniformly to one.  This is
Corollary~\ref{cor:profile-zero-uniform}; it follows from vanishing
affine variance and the uniform reference exponential moments.
The tilt below is fixed before the limit in \(k\).

\begin{proposition}\label{prop:mass-polynomial}
Fix \(w'>T\), positive integers \(r_k\) satisfying
\(\delta_k\log_Lr_k\to w'\), and constants \(0<c_0<C_0<\infty\).
For every \(K>0\), all sufficiently large \(k\) satisfy
\begin{equation}\label{eq:mass-polynomial}
 \sup_{c_0r_k\le\|x\|_\infty\le C_0r_k}C_{b_k}(x)
                                      \le r_k^{-K}.
\end{equation}
\end{proposition}

\begin{proof}
Choose \(T<u<v<w'\), and use radii
\[
 \mathcal R_k=\{4^j:j\in\Z,\quad
                         u\le\delta_k\log_L(4^j)\le v\}.
\]
Their number \(N_k\) obeys
\begin{equation}\label{eq:mass-annulus-count}
 N_k=\frac{v-u}{\delta_k\log_L4}+O(1),\qquad
 \frac{N_k}{\log r_k}\longrightarrow
                          q:=\frac{v-u}{w'\log4}>0.
\end{equation}
For \(R\in\mathcal R_k\), the annuli
\(R<|z|<2R\) about the source at \(0\) are disjoint.
Their maximal radius is \(o(r_k)\); therefore they avoid the other
source \(x\), uniformly over the displacements in the proposition.
Figure~\ref{fig:mass-geometry} compares this annular geometry with
the pinned geometry used for the lower bound.
\begin{figure}[tb]
\centering
\begin{tikzpicture}[>=Latex,font=\small]
\begin{scope}[xshift=2.3cm,yshift=2.3cm]
 \fill[gray!18] (-2,-2) rectangle (2,2);
 \fill[white] (-1.5,-1.5) rectangle (1.5,1.5);
 \draw[gray!70] (-2,-2) rectangle (2,2);
 \draw[gray!70] (-1.5,-1.5) rectangle (1.5,1.5);
 \node at (0,1.75) {$h=0$};
 \node[anchor=north east] at (1.48,1.48) {$D$};
 \fill[blue!18] (-1.03,-.38) rectangle (-.53,.38);
 \fill[orange!22] (.53,-.38) rectangle (1.03,.38);
 \node at (-.78,0) {$f$};
 \node at (.78,0) {$g$};
 \draw[blue!65!black,thick] (0,0) ellipse[x radius=1.27cm,y radius=.82cm];
 \draw[->,blue!65!black,thick] (-.36,.786)--(-.1,.818);
 \draw[<->] (-.53,-1.13)--(.53,-1.13) node[midway,below] {$c n$};
 \node[align=center] at (0,-2.5) {Pinned covariance below $T$};
\end{scope}
\begin{scope}[xshift=8.6cm,yshift=2.3cm]
 \fill[blue!12,even odd rule] (0,0) circle (1.68) (0,0) circle (1.4);
 \fill[blue!12,even odd rule] (0,0) circle (.42) (0,0) circle (.35);
 \draw[blue!55!black] (0,0) circle (1.4);
 \draw[blue!55!black] (0,0) circle (1.68);
 \draw[blue!55!black] (0,0) circle (.35);
 \draw[blue!55!black] (0,0) circle (.42);
 \fill (0,0) circle (.035) node[below left] {$0$};
 \fill (2.6,0) circle (.035) node[right] {$x$};
 \draw[dashed,gray] (.45,0)--(2.53,0);
 \node at (.9,1.05) {$4R$};
 \node[below right,font=\footnotesize] at (.17,-.42) {$R$};
 \draw[->,blue!60!black] (-.33,.86) arc[start angle=111,end angle=145,radius=.92];
 \node[align=center] at (0,-2.08) {$\delta\log_LR_i\in[u,v]\subset(T,w')$};
 \node[align=center] at (0,-2.5) {Circulation tilts above $T$};
\end{scope}
\end{tikzpicture}
\caption{Two finite geometries for the mass comparison. Left: a primary pin region of positive width surrounds $D$. A loop contributing to the covariance of two separated density smears must have diameter at least their separation. The displayed loop is schematic. Right: geometrically spaced annuli surround one spin insertion and remain inside the separation scale of the other. Only two annuli are drawn. Their logarithmic scales lie in an interval where the height coefficient vanishes, so a fixed circulation tilt has arbitrarily small fluctuation cost on each annulus.}
\label{fig:mass-geometry}
\end{figure}


Choose a nonnegative, nonzero smooth radial function \(\chi\)
supported in \((1,2)\), and let
\[
       p(z)=\chi(|z|)\nabla\arg z,
       \qquad
       D=\int_{\R^2}p(z)\cdot\nabla\arg z\,dz>0.
\]
This is a smooth divergence-free form, compactly supported in the
unit annulus.  Choose its radial stream function \(\psi\), constant
on each complementary component, with \(\nabla^\perp\psi=p\).
At scale \(R\), sample \(\psi_R(z)=\psi(z/R)\) on the faces of
the primary height graph and take its rotated discrete difference.
Denote the resulting edge form by \(\eta_R\), choosing its sign
to agree with the positive source.  Thus, on an edge with midpoint
\(z_e\), its corresponding coordinate is
\(R^{-1}p(z_e/R)+O(R^{-2})\).
It has zero discrete divergence,
\begin{equation}\label{eq:mass-annulus-form}
 \max_e|\eta_R(e)|\le C/R,\qquad
 \sum_e|\eta_R(e)|^2\le C.
\end{equation}
For any affine height difference \(\nabla h+\alpha\) representing
the unit spin twist at \(0,x\), discrete summation by parts gives
\begin{equation}\label{eq:mass-deterministic-circulation}
 \sum_e\eta_R(e)(\nabla h+\alpha)_e=D_R,
                       \qquad D_R\longrightarrow D>0.
\end{equation}
Indeed the gradient part pairs to zero.  The remaining pairing is
the difference of the stream-function values at the two source
faces, multiplied by their circulations \(\pm2\pi\).
The stream function is constant near each source, so there is no
dependence on the height configuration or on the path chosen for
\(\alpha\).  The convergence in
\eqref{eq:mass-deterministic-circulation} also follows directly by
Riemann sums; its error is uniform here because the second source
lies outside every annulus.  Fix \(d_{\rm circ}>0\) so that \(D_R\ge d_{\rm circ}\)
for all sufficiently large radii.

We make the partition-ratio estimate at fixed \(k,x\) in a finite
free spin box large enough to contain all these annuli in its
interior.  Write \(Z_\alpha\) and \(Z_0\) for the twisted and
centered height sums, both with exterior face height zero.
Free-boundary duality gives \(C_{G,b_k}(0,x)=Z_\alpha/Z_0\).
For any fixed \(\lambda>0\), insertion of the deterministic tilt
from \eqref{eq:mass-deterministic-circulation} gives the exact identity
\begin{equation}\label{eq:mass-tilted-ratio}
 e^{\lambda\sum_{R\in\mathcal R_k}D_R}C_{G,b_k}(0,x)
 =\frac{Z_\alpha\left(
       e^{\lambda\sum_R(\eta_R,\nabla h+\alpha)}\right)}{Z_0}.
\end{equation}
Here \(Z_\alpha(F)\) means the unnormalized twisted sum with
insertion \(F\).  We will bound the right side independently of
the outer spin box.

Fix a small relative cell size \(t>0\).  Cover the support of each
\(\eta_R\) by square primary cells of side
\(m_R=\lfloor tR\rfloor\), with disjoint vertex sets, leaving
individual edges between cells as seam increments.  The support
of \(\chi\) has positive distance from the annular boundary.
Thus for sufficiently small fixed \(t\), every cell meeting this
support lies inside the annulus, avoids both sources, and is
simply connected.  There are \(O(t^{-2})\) such cells at each
radius.  Edges not in a retained cell and carrying a nonzero tilt
number \(O_t(R)\).

Release the endpoint constraints between these cells and the seam
increments, and release all remaining constraints outside the
cells.  The tilted affine-to-centered comparison
\cite[Lemma~2.2(ii)]{X} bounds the normalized expression in
\eqref{eq:mass-tilted-ratio} by the product of the cell transforms
and the independent whole-edge seam transforms.  Untilted factors
on the rest of the graph have affine-to-centered ratio at most
one.  On a retained cell the connection is an integral gradient,
\(\alpha=\nabla\kappa\); reindexing \(h'=h+\kappa\) removes it
from both the weight and the statistic being tilted.  Accordingly
the factor for that cell is the centered free transform
\[
             \E_{Q,b_k}^{\rm free}
                     \exp\{\lambda(\eta_R,\nabla h)_Q\}.
\]
In particular, this gauge change introduces no additional real
exponential factor.  Released constant modes are taken modulo
common translations; the edge-difference tests annihilate every
such constant.  Equivalently, one can retain vanishing quadratic
penalties on these constants in both sectors and cancel their
identical factors before removing them.  No outer-box normalizer
is present in the resulting bound.

We now estimate the factors for one annulus.  The total squared
coefficient cost of its seam increments is \(O_t(R^{-1})\) by
\eqref{eq:mass-annulus-form}.  The Bessel increment transform is
\(\exp\{b_k(\cosh(2\pi s)-1)\}\), so their combined logarithmic
cost is \(O_{t,\lambda}(R^{-1})\), uniformly in \(k\).
On each retained square replace \(\eta_R\) by its value at a
chosen point of the square.  Smoothness of \(p\) makes the total
squared cost of the coefficient error over all cells
\(O(t^2)+o_R(1)\).  The bounded-flow estimate \eqref{eq:bounded-flow-mgf} and
H\"older's inequality with exponent two therefore give
\begin{align}
 &\sum_Q\log\E_{Q,b_k}^{\rm free}
                  e^{\lambda(\eta_R,\nabla h)_Q}
 \notag\\[-2pt]
 &\qquad\le
 \frac12\sum_Q\log\E_{Q,b_k}^{\rm free}
                    e^{2\lambda(\bar\eta_{R,Q},\nabla h)_Q}
                       +C_\lambda t^2+o_{t,\lambda,R}(1),
 \label{eq:mass-affine-cell-reduction}
\end{align}
where \(\bar\eta_{R,Q}\) is the constant coefficient on \(Q\).
The error coefficients multiplied by the fixed tilt are uniformly
bounded, as required by \eqref{eq:bounded-flow-mgf}.  That estimate
bounds each error transform by its squared coefficient cost.
Summing the costs gives the total \(C_\lambda t^2\) in
\eqref{eq:mass-affine-cell-reduction}.

Each constant test in the sum is a linear combination of the two
square affine observations defining \(a_{m_R}(b_k)\).  Its
coefficient vector has size at most \(C_\lambda m_R/R\le C_\lambda t\).
Translations of the square do not change its law.  Moreover,
\[
       \delta_k\log_Lm_R
            =\delta_k\log_LR+\delta_k\log_Lt+o(1).
\]
Thus these scales lie in a slightly enlarged compact subinterval
of \((T,w')\) for all large \(k\).  By
Corollary~\ref{cor:profile-zero-uniform}, the transforms of these
affine tests tend uniformly to one, uniformly also in their
bounded coefficient vectors.  Since the number of cells is uniformly bounded
once \(t\) is fixed, their total logarithmic cost is \(o_k(1)\)
uniformly over \(R\in\mathcal R_k\).
The reference estimates and \(\min\mathcal R_k\to\infty\) also
make all seam and discretization errors uniform.  We obtain,
for every fixed \(\lambda,t\), the per-annulus bound
\begin{equation}\label{eq:mass-annulus-cost}
           \log(\text{product of its transforms})
                             \le C_\lambda t^2+o_k(1).
\end{equation}

Choose \(\lambda\) first, and then \(t\) so small that
\(C_\lambda t^2\le\lambda d_{\rm circ}/4\).  For all sufficiently large
\(k\), the uniform error in \eqref{eq:mass-annulus-cost} is at most
\(\lambda d_{\rm circ}/4\).  Equations~\eqref{eq:mass-tilted-ratio} and
\eqref{eq:mass-annulus-count} imply
\[
                    C_{G,b_k}(0,x)
                         \le\exp\{-\lambda d_{\rm circ}N_k/2\}.
\]
This estimate holds uniformly in every sufficiently large outer
spin box, at the fixed values of \(b_k,r_k,x\).  We may therefore
first take the prescribed free-box thermodynamic limit.  The
resulting estimate for \(C_{b_k}(x)\) is uniform in the indicated
range of \(x\).  Finally choose the fixed \(\lambda\) large enough
that \(\lambda d_{\rm circ}q/4>K\).  Equation~\eqref{eq:mass-annulus-count}
then gives \eqref{eq:mass-polynomial}.
\end{proof}

\subsection{A finite-size criterion and completion of the proof}

The Lieb--Rivasseau inequality for plane rotators
\cite{Lieb,Rivasseau}, in the form stated in
\cite[Lemma~17]{vEL}, says that for a finite subgraph \(H\subset G\),
\(u\in H\), and \(v\notin H\),
\begin{equation}\label{eq:mass-lieb-rivasseau}
 C_{G,b}(u,v)\le\sum_{z\in\partial H}
                       C_{H,b}(u,z)C_{G,b}(z,v).
\end{equation}
Here \(\partial H\) is the set of vertices of \(H\) adjacent to a
vertex outside \(H\); the first factor is the correlation with
free boundary on \(H\).  We give the iteration so its conclusion
refers to exactly the axis mass of this paper.

\begin{lemma}\label{lem:mass-finite-size}
Let \(H_r=[-r,r]^2\cap\Z^2\), and put
\[
                 S_{r,b}=\sum_{z\in\partial H_r}C_{H_r,b}(0,z).
\]
If \(S_{r,b}\le1/2\), then \(m(b)\ge(\log2)/r\).
\end{lemma}

\begin{proof}
At fixed \(b,r\), take the free-box thermodynamic limit in
\eqref{eq:mass-lieb-rivasseau}.  For a target \(se_1\), apply the
inequality successively in translates of \(H_r\) centered at the
boundary point reached in the previous step.  Each step moves the
center by at most \(r\) in the supremum norm.  Hence at least
\(\lfloor s/r\rfloor-1\) iterations are possible before the target
can belong to the next translated square.  Translation invariance
and \(C_b\le1\) give
\[
                   C_b(se_1)\le
                    S_{r,b}^{\lfloor s/r\rfloor-1}.
\]
Taking \(-s^{-1}\log\) and then \(s\to\infty\) proves the claim.
\end{proof}

\begin{proposition}\label{prop:mass-upper-length}
Along the extracted sequence,
\begin{equation}\label{eq:mass-limsup}
                 \limsup_k\delta_k\log_L\frac1{m(b_k)}\le T.
\end{equation}
\end{proposition}

\begin{proof}
Fix \(w'>T\), and set \(r_k=\lfloor L^{w'/\delta_k}\rfloor\).
Proposition~\ref{prop:mass-polynomial}, with any fixed \(K>1\),
and finite-to-infinite monotonicity imply
\[
 S_{r_k,b_k}\le\sum_{z\in\partial H_{r_k}}C_{b_k}(z)
                     \le8r_k\,r_k^{-K}<\frac12
\]
for all sufficiently large \(k\).  Thus
Lemma~\ref{lem:mass-finite-size} yields
\(m(b_k)\ge(\log2)/r_k\).  Consequently
\[
 \limsup_k\delta_k\log_L\frac1{m(b_k)}
 \le\lim_k\delta_k\log_L\frac{r_k}{\log2}=w'.
\]
Letting \(w'\downarrow T\) proves \eqref{eq:mass-limsup}.
\end{proof}

\begin{proof}[Proof of Theorem~\ref{thm:main}]
Start with any sequence \(b_k<b_c\) tending to \(b_c\).
The construction in Sections~\ref{sec:trajectory} and
\ref{sec:profile} admits a subsequence with the profile
\eqref{eq:mass-profile}.  Its transition point is
\[
                       T=\frac{\pi}{\sqrt{3\mathcal P}},
\]
where \(\mathcal P\in(0,\infty)\) is independent of the extracted
sequence.  Propositions~\ref{prop:mass-lower-length} and
\ref{prop:mass-upper-length} prove
\eqref{eq:mass-scale-goal} on that subsequence.
If the asserted limit failed for the original parameter approach,
there would be a sequence on which its distance from \(T\) stayed
bounded below; applying the same extraction to that sequence
would give a contradiction.  Therefore the logarithmic-scale
limit holds as \(b\uparrow b_c\) without extraction.
Converting from base \(L\) to natural logarithms gives
\[
 \lim_{b\uparrow b_c}\sqrt{b_c-b}\log\frac1{m(b)}
       =(\log L)\,T
       =\frac{\pi\log L}{\sqrt{3\mathcal P}}
       =:A_{\rm XY}\in(0,\infty).
\]
All spin correlations used above are limits of the prescribed
free spin boxes before their separation tends to infinity.  Thus
the mass in this identity is exactly the mass in
Theorem~\ref{thm:main}.
\end{proof}
